\section{Dated Timeline}

The final column distinguishes a recorded event from the later witness that reports it. Approximate ranges and disputed dates remain explicit.

\begin{longtable}{>{\bfseries\RaggedRight\arraybackslash}p{0.16\linewidth}>{\RaggedRight\arraybackslash}p{0.48\linewidth}>{\RaggedRight\arraybackslash}p{0.27\linewidth}}
\toprule
\textbf{Date} & \textbf{Event or development} & \textbf{Witness and boundary}\\
\midrule
\endfirsthead
\toprule
\textbf{Date} & \textbf{Event or development} & \textbf{Witness and boundary}\\
\midrule
\endhead
\bottomrule
\endfoot

1840s--60s & Polish-speaking migrants, many from Prussian-ruled Polish lands, enter Milwaukee's industrial and service economy; South Side settlement grows. & Archdiocese, UWM, National Register, and Portal Polonii syntheses; regional origins were plural.\\

1863 & Organizing of a Polish Catholic congregation is reported as begun under Father Buczinski. & Gregory's 1931 city history, p.~902; later official histories use 1866 for formal foundation.\\

1866 & Thirty families formally organize Saint Stanislaus, acquire a former Lutheran brick church at Fifth and Mineral, and establish Milwaukee's first Polish parish. & Official, preservation, and later historical accounts converge; no founding minutes were checked.\\

1866--78 & Polish parish instruction begins and receives dedicated school space. & Sources attach 1866, 1867, 1872, or 1873 to instruction and 1873 or 1878 to a schoolhouse; exact sequence unresolved.\\

1871 & Saint Hedwig becomes Milwaukee's second Polish parish, serving a constituency for whom Saint Stanislaus was distant. & Gregory, UWM, and Portal Polonii; first clear daughter-parish phase.\\

July 1872 & Cornerstone laid for the present Saint Stanislaus Church at Fifth and Mitchell. & HABS, AHI, Gregory, National Register, and parish history.\\

1873 & A portion of the present church is dedicated and placed in use. & Gregory and institutional memory; principal construction continues in 1884--94.\\

1870s & Kashubian fishing settlement develops on Jones Island, with many families connected sacramentally to Saint Stanislaus. & Archdiocesan and Milwaukee Polonia histories; island and parish foundations are not identical events.\\

1882 & Parish history reports 96 marriages, 589 baptisms, and 520 confirmations in the year. & Institutional retrospective; definitions and original registers were not audited.\\

1883 & Saint Hyacinth is organized as Polish settlement expands on the South Side. & UWM and contextual parish histories; evidence of Saint Stanislaus's generative rather than merely centralized role.\\

1884--94 & Principal later construction completes sanctuary, copper roof, organ, and painted interior. & HABS physical history; demonstrates why 1873 is a use date, not a universal completion date.\\

1885 & Henry Messmer designs \$17,000 in renovations. & Wisconsin AHI HI27246; digital record gives value without full schedule.\\

4--5 May 1886 & Polish workers gather at Saint Stanislaus before marching toward the Bay View rolling mill; militia fire the next day kills seven people. & John Gurda account hosted by Wisconsin Labor History Society; gathering place does not establish parish sponsorship of the strike.\\

1888 & Saint Vincent de Paul and Saint Josaphat are among the Polish congregations formed as settlement spread. & UWM and Polish-parish histories; daughter institutions had their own founders and later histories.\\

1894--95 & Four great bells and successive marble work mark a major interior and tower campaign. & Detailed parish history; one 2018 feature dates bells to 1882, so installation year remains source-dependent.\\

1895--1913 & Entrances, tower windows, buttresses and apostle statues, crossing flèche, marble altars, pulpit, and decoration reshape church exterior and interior. & HABS; individual years within the range cannot all be recovered.\\

1914--15 & Auxiliary Bishop Edward Kozłowski is associated with Saint Stanislaus during his brief Milwaukee episcopate. & Later Polish institutional history; a full episcopal-residence and event chronology was not reconstructed.\\

1926 & Present principal school front built to plans by Herbst and Kuenzli dated 16 February. & Wisconsin AHI HI118139 and National Register nomination; controls over an ``early 1930s'' parish recollection.\\

1931 & John G. Gregory publishes a city history preserving the 1863 organization, reused first church, 1873 dedication, and 1878 schoolhouse chronology. & Near-contemporary synthesis, not a founding document.\\

1932 & Saint Stanislaus High School opens with seventy-two freshmen, boys and girls, and grows into a regional South Side Catholic school. & A 1946 \emph{Milwaukee Sentinel} retrospective preserved by Marquette University; school-program history is distinct from the 1926 building date.\\

1941 & A parish diamond-jubilee memorial volume is produced. & HABS cites the volume; it was not examined in full here.\\

1947 & The state inventory records the school's name change to Notre Dame High School. & Wisconsin AHI HI118139; the building had already housed secondary instruction since 1932.\\

1950s--60s & Expressway clearance and postwar suburban mobility disrupt South Side residence and accelerate Polish dispersal from the old parish neighborhood. & National Register context, Portal Polonii, and later reporting; no annual parish census was found.\\

1958 & Raymond A. Punda becomes pastor and prepares an extensive centennial renovation. & Institutional history.\\

1960 & Towers and domes rebuilt, gilded, and reset; clocks and masonry renewed. & Detailed institutional chronology; other inventories use broader early-1960s or 1966 dating.\\

1962 & Fire damages a side altar and smokes the church interior. & Institutional history; no fire report checked.\\

1962/64 & Old stained glass is removed and dalle de verre installed. & The City of Milwaukee survey says 1962; institutional history says 1964; no checked permit or contract resolves the difference.\\

1963--64 & Rectory remade; ambulatory and chapel join it to church. & Parish and preservation records; remodeling-versus-replacement language remains unsettled.\\

1965--66 & Travertine, freestanding altar, relocated pulpit and font, reduced rail, extended sanctuary, and Częstochowa mosaic complete the centennial-era remaking. & Institutional history; the campaign combined maintenance, safety, commemorative, artistic, and liturgical judgments.\\

1966 & A 31-page parish centennial history is produced. & Notre Dame Archives catalogs the volume; it was not examined in full here.\\

1970s--80s & Parish numbers decline and the building's future becomes uncertain. & Institutional history and later journalism; no complete quantitative series located.\\

Reported 1988 & Notre Dame High School closes; the building later remains in consolidated Catholic elementary-school use. & OnMilwaukee retrospective and National Register nomination; no school-board or archdiocesan closing act checked.\\

Reported 2003 & Sale of the church is described as under archdiocesan consideration. & 2014 OnMilwaukee oral-history feature; not supported here by a contemporaneous sale file.\\

15 October 2007 & ICKSP announces Archbishop Timothy Dolan's invitation to serve the traditional-Latin-Mass apostolate in Milwaukee, with Sunday Mass at Saint Stanislaus to begin in Advent. & Contemporary institutional announcement; it records invitation and planned ministry, not transfer of title.\\

2007--08 & Parish placed in the Institute's care; institutional history dates erection of Saint Stanislaus Oratory to 2008; restoration begins. & ICKSP history. A 2018 feature instead gives 2016 for canonical establishment; no decree located.\\

2008 onward & Sanctuary and high altar reconstructed, floor and rail renewed, and historic liturgical furnishings recovered in phases. & ICKSP and \emph{Catholic Herald}; no single completion date.\\

2011--12 & Organ restored. & Institutional history; scope detailed in building ledger.\\

2013--14 & Sacristy and former connecting chapel restored and adapted for current service. & Institutional history.\\

2014/15--21 & Reconstructed stained-glass windows installed in phases using historical evidence and surviving fragments. & Parish, contractor, and press records disagree slightly on discovery date; completion reported in 2021.\\

2016 & Interior decorative and marble restoration advances; 24--25 September sesquicentennial observances include Benediction with Archbishop Jerome Listecki and a pontifical Mass with Auxiliary Bishop Joseph Perry, who consecrates the rebuilt main altar. & Contemporary ICKSP anniversary account; parish anniversary, not church-building sesquicentennial or oratory-erection proof.\\

2017--18 & Choir loft redesigned and renewed. & Institutional history.\\

9 November 2018 & West Mitchell Street Commercial Historic District entered in the National Register; church, school, and rectory are contributing resources. & Wisconsin AHI and National Register record. District contribution is not an individual National Register listing for the church.\\

December 2018--December 2019 & Copper-dome and roof campaign proceeds with masonry, doors, and site restoration. & Professional roofing account and institutional history.\\

2019--21 & Chandeliers, rectory interiors, and final stained-glass phases complete major visible portions of the program. & Institutional history; ongoing maintenance means restoration has no final terminal date.\\

19 November 2022 & Saint Stanislaus observes the 150th anniversary of the present church's 1872 building campaign. & ICKSP event notice; distinct from the parish's 150th anniversary in 2016.\\

July--August 2023 & All four great bells receive replacement clappers after scheduled maintenance finds the original striking ends worn flat; the original clappers are retained for display. & Contemporary parish bulletin; maintenance of working bells, not replacement of the bells.\\

28 October 2023 & Auxiliary Bishop Jeffrey Haines confirms seventy-eight young faithful at the oratory. & Contemporary ICKSP report; a sacramental event and activity snapshot, not a resident-parish census.\\

17 March 2024 & Parish and oratory announce planned acquisition and renovation of two nearby buildings for classrooms, meetings, and parking. & Institutional announcement; later public records establish purchase and approvals, but not completed renovation.\\

20 September 2024 & Love One Another page reports \$11,705 pledged against a \$36,183 parish goal and names irrigation and exterior lighting needs. & Dated archdiocesan campaign snapshot; not a total restoration cost.\\

17 October 2024 & St Stanislaus Properties LLC purchases 501--505 West Historic Mitchell Street for \$875,000. & City assessor transaction record; proves acquisition, not finished rehabilitation or occupancy.\\

7 July / 11 September 2025 & Historic Preservation Commission approves replacement brick facing; Board of Zoning Appeals grants special use as a religious assembly hall subject to permits and occupancy certificates. & Official city records. These are rehabilitation and use approvals, not a completion certificate.\\

3 May / 4 June 2026 & Twenty-five children receive first Communion; hundreds participate in Corpus Christi Mass and a five-block procession to Saint Anthony. & Contemporary ICKSP event reports; evidence of recurrent use and regional community, not proof of neighborhood demographic reversal.\\

19 July 2026 & Historical and mutable-status cutoff. & Current schedule, staff, public status, campaign, and expansion evidence checked to this date; independent specialist review remains outstanding.\\
\end{longtable}
